\section{An explicit economy with heterogeneous households}\label{sec:corner}

\subsection{Corner economies}

Let the employed save up to the cap and the unemployed spend everything. Under the employment
design of \citet{krusellsmith1998} --- the employment rate at shock $z$ is $m(z)$, consistent under
the transitions --- the cross-section after yesterday's shock $y$ and today's $z$ puts mass on two
asset levels, $\abar$ for those employed last period and $0$ for the others, jointly with the current
employment state. Capital is $\abar\, m(y)$. Below the root of a history's tree, capital and labour
depend only on the shocks along the node and today's shock, so yesterday's shock enters prices only
through today's capital: $(z,K)$-sufficiency holds with no injectivity.

The corner policy is optimal when three conditions hold along the price path: labour income covers
the floor plus the corner saving; for the employed, $u'(c)\le\beta\E[R'u'(c')]$ at the cap; for the
unemployed, $u'(c)\ge\beta\E[R'u'(c')]$ at zero, where $c'$ is consumption next period under the
corner policy itself. The proof is a verification argument: the value of following the corner policy
is today's utility plus a term independent of assets; against it the objective is concave, the
tangent lines today and tomorrow bound the gain from leaving the corner by
$(g-a)\,(u'(c)-\beta\E R'u'(c'))$, which has the right sign; so the policy's value solves the Bellman
equation and is the value function. With prices bounded in $[R_{lo},R_{hi}]\times[w_{lo},w_{hi}]$ and
marginal utility decreasing, the conditions hold along every price path once they hold at their
worst cases.
\formal{policy\_eq\_cornerPolicy, policyValue\_eq\_valueFunction, cornerConditions\_of\_bounds}

\subsection{The economy}

\begin{center}
\begin{tabular}{ll}
\toprule
Aggregate shocks & two, i.i.d.\ with probability $1/2$ \\
Employment & efficiency $2$ (employed), $1/4$ (unemployed) \\
Employment rates & $4/5$ (bad), $17/20$ (good); job finding $49/50$ \\
Job loss & $49/200,\ 73/400,\ 197/850,\ 147/850$ by (yesterday, today) \\
Preferences & logarithmic, $\beta=24/25$; asset cap $\abar=1/16$ \\
Technology & Cobb--Douglas, $\alpha=1/10$, TFP $1/4$, $\delta=1/10$ \\
Boxes & $K^{1/10}\in[11/15,\,3/4]$, $L^{1/10}\in[21/20,\,16/15]$ \\
\bottomrule
\end{tabular}
\end{center}

\noindent
The box endpoints are tenth powers of rationals, so the powers $1/10$ and $9/10$ in the prices are
exact and every inequality is checked in exact rational arithmetic. Gross returns lie in about
$[1.42,1.63]$ and wages in about $[0.155,0.161]$. Capital is $1/20$ after a bad shock and $17/320$
after a good one, labour $33/20$ or $139/80$: both inside the boxes at every node of the equilibrium
price tree, so prices are the true Cobb--Douglas marginal products.

\begin{theorem}\label{thm:corner}
In this economy: sequential equilibrium is unique from every initial state; the sequence-state
equilibrium exists, is unique, and is the corner one; $(z,K)$-sufficiency holds and a recursive
equilibrium on $(z,K)$ represents it; and its law of motion on $(z,K)$ has a unique stationary
distribution on the capital sets --- weight $1/4$ on each pair of today's shock and capital from
yesterday's shock --- which is ergodic.
\end{theorem}
\formal{cornerUnique, cornerWitness, cornerStationary, ergodic\_statLaw, seqPrices\_mem\_box}

\noindent
Uniqueness comes cheaply because the corner conditions hold along every price path: the household's
policy is the corner one whatever it expects, so the distribution path does not depend on prices.
Three features limit the example. Assets sit at corners rather than at interior precautionary
levels. The gross return exceeds $1/\beta$, which is what makes the cap bind for the employed. And
with capital share $1/2$ and labour inside its box the corner economy cannot exist: the employed's
net income over the unemployed's spending is bounded by roughly $(1-\alpha)/\alpha$, too small for
both corner conditions.
